\documentclass[11pt,oneside,openany]{book}
\usepackage[letterpaper,left=.9in,right=.9in,top=.9in,bottom=.9in,headheight=16pt,headsep=.22in]{geometry}
\usepackage{fontspec}
\setmainfont{P052-Roman.otf}[Path=fonts/,BoldFont=P052-Bold.otf,ItalicFont=P052-Italic.otf,BoldItalicFont=P052-BoldItalic.otf]
\setsansfont{NimbusSans-Regular.otf}[Path=fonts/,Scale=.94,BoldFont=NimbusSans-Bold.otf,ItalicFont=NimbusSans-Italic.otf,BoldItalicFont=NimbusSans-BoldItalic.otf]
\setmonofont{DejaVu Sans Mono}[Scale=.76]
\usepackage{amsmath,amssymb,mathtools,amsthm,bm,microtype}
\usepackage{graphicx,float,booktabs,longtable,tabularx,array,multirow,ragged2e}
\usepackage[table,dvipsnames]{xcolor}
\usepackage{enumitem,fancyhdr,titlesec,xurl,needspace,calc,fvextra,listings,makeidx}
\usepackage[most]{tcolorbox}
\usepackage{hyperref,bookmark}
\usepackage[nameinlink,capitalize,noabbrev]{cleveref}
\definecolor{Ink}{HTML}{183247}
\definecolor{Teal}{HTML}{39766E}
\definecolor{Muted}{HTML}{596773}
\definecolor{Pale}{HTML}{EFF4F4}
\definecolor{Coral}{HTML}{B66151}
\definecolor{Gold}{HTML}{A98440}
\colorlet{ink}{Ink}\colorlet{teal}{Teal}\colorlet{accent}{Teal}\colorlet{muted}{Muted}\colorlet{pale}{Pale}\colorlet{light}{Pale}\colorlet{red}{Coral}\colorlet{gold}{Gold}
\colorlet{deepblue}{Ink}\colorlet{rose}{Coral}\colorlet{restrictedblue}{Ink}\colorlet{ambientrose}{Coral}
\colorlet{proofgreen}{Pale}\colorlet{recordblue}{Pale}\colorlet{derivedgold}{Gold!8}\colorlet{openred}{Coral!7}
\colorlet{provedgreen}{Pale}\colorlet{computedblue}{Pale}\colorlet{skeletonyellow}{Gold!8}\colorlet{conjecturered}{Coral!7}\colorlet{rulegray}{Muted!30}
\hypersetup{colorlinks=true,linkcolor=Ink,citecolor=Teal,urlcolor=Teal,pdftitle={Realization Calculus: Quantum and Circuits},pdfauthor={Komorebi P.},pdfsubject={Arithmetic, Symmetry, and the Cost of Realization}}
\urlstyle{tt}\raggedbottom
\setlength{\parindent}{0pt}\setlength{\parskip}{5pt plus 1pt minus 1pt}
\setlength{\emergencystretch}{4em}
\setlist{nosep,leftmargin=1.6em,topsep=4pt}
\setcounter{tocdepth}{0}\setcounter{secnumdepth}{2}
\numberwithin{equation}{chapter}
\let\cleardoublepage\clearpage
\pagestyle{fancy}\fancyhf{}
\fancyhead[L]{\small\sffamily\color{Muted}REALIZATION CALCULUS}
\fancyhead[R]{\small\sffamily\color{Muted}Quantum and Circuits}
\fancyfoot[R]{\small\sffamily\thepage}
\renewcommand{\headrulewidth}{.3pt}
\fancypagestyle{plain}{\fancyhf{}\fancyfoot[R]{\small\sffamily\thepage}\renewcommand{\headrulewidth}{0pt}}
\titleformat{\chapter}[display]{\normalfont\sffamily\bfseries\color{Ink}\raggedright\hyphenpenalty=10000}{\small\MakeUppercase{\chaptertitlename}\ \thechapter}{12pt}{\fontsize{26}{31}\selectfont}
\titlespacing*{\chapter}{0pt}{0pt}{22pt}
\titleformat{\section}{\large\sffamily\bfseries\color{Ink}}{\thesection}{.65em}{}
\titleformat{\subsection}{\normalsize\sffamily\bfseries\color{Ink}}{\thesubsection}{.65em}{}
\titleformat{\subsubsection}{\normalsize\sffamily\bfseries\color{Ink}}{}{0pt}{}
\newtheoremstyle{volume}{8pt}{5pt}{\itshape}{}{\bfseries\color{Ink}}{.}{.6em}{}
\theoremstyle{volume}
\newtheorem{theorem}{Theorem}[chapter]
\newtheorem{proposition}[theorem]{Proposition}\newtheorem{lemma}[theorem]{Lemma}\newtheorem{corollary}[theorem]{Corollary}\newtheorem{principle}[theorem]{Principle}\newtheorem{counterexample}[theorem]{Counterexample}
\theoremstyle{definition}\newtheorem{definition}[theorem]{Definition}\newtheorem{example}[theorem]{Example}\newtheorem{algorithm}[theorem]{Algorithm}\newtheorem{protocol}[theorem]{Protocol}
\theoremstyle{remark}\newtheorem{remark}[theorem]{Remark}\newtheorem{question}[theorem]{Question}\newtheorem{conjecture}[theorem]{Conjecture}
\tcbset{enhanced,breakable,arc=0pt,boxrule=.3pt,colframe=Muted!35,left=9pt,right=9pt,top=7pt,bottom=7pt,fonttitle=\small\sffamily\bfseries,coltitle=Ink,colbacktitle=Pale}
\newtcolorbox{provedbox}[1][]{colback=Pale,title={Proved statement},#1}
\newtcolorbox{recordbox}[1][]{colback=Pale,title={Recorded computation},#1}
\newtcolorbox{derivedbox}[1][]{colback=Gold!8,title={Derived theorem or corollary},#1}
\newtcolorbox{frontierbox}[1][]{colback=Coral!6,title={Open program or boundary},#1}
\newtcolorbox{scopebox}[1][]{colback=Pale,title={Scope boundary},#1}
\newtcolorbox{computedbox}[1][]{colback=Pale,title={Exact finite computation},#1}
\newtcolorbox{skeletonbox}[1][]{colback=Gold!8,title={Research template},#1}
\newtcolorbox{conjecturebox}[1][]{colback=Coral!6,title={Conjectural frontier},#1}
\newtcolorbox{boundarybox}[1][]{colback=Pale,title={Scope boundary},#1}
\newtcolorbox{circuitscopebox}[1]{colback=Pale,title={#1}}
\newenvironment{statusbox}[1]{\begin{circuitscopebox}{#1}}{\end{circuitscopebox}}
\newcolumntype{L}[1]{>{\raggedright\arraybackslash}p{#1}}
\providecommand{\tightlist}{\setlength{\itemsep}{0pt}\setlength{\parskip}{0pt}}
\newcommand{\lead}[1]{{\fontsize{12}{16}\selectfont #1}\par}
\newcommand{\note}[1]{\begin{tcolorbox}[colback=Pale]#1\end{tcolorbox}}
\newcommand{\fig}[3]{\begin{center}\includegraphics[width=#2\linewidth,height=.46\textheight,keepaspectratio]{assets/q_#1.pdf}\end{center}{\small\color{Muted}#3}\par}
\newcommand{\guidepage}[1]{\clearpage\section*{#1}}
\newcommand{\evidence}[1]{\par{\small\color{Muted}#1}\par}
\newcommand{\sourcepath}[1]{}
\newcommand{\visualfigure}[2]{\begin{center}\includegraphics[width=\textwidth,height=3.15in,keepaspectratio]{assets/c_#1.pdf}\par\vspace{3pt}{\small #2}\end{center}}
\newcommand{\circuitscope}[1]{\begin{quote}\small\textbf{Scope.} #1\end{quote}}
\newcommand{\source}[1]{}
\newcommand{\claim}[1]{\par{\small\sffamily\bfseries\color{Teal}CLAIM #1}\par\medskip}
\newcommand{\pathref}[1]{{\small\texttt{\detokenize{#1}}}}
\newcommand{\code}[1]{{\small\texttt{\detokenize{#1}}}}
\newcommand{\resultrow}[4]{\ref{#1} & #2 & #3 & \pageref{#1}\\}
\newcommand{\pagehead}[2]{\clearpage\section*{#1}\addcontentsline{toc}{section}{#1}{\small\color{Teal}\sffamily #2}\par\medskip}
\newcommand{\TC}{\operatorname{tc}}\newcommand{\Deg}{\operatorname{DegreeLE}}\newcommand{\x}{\mathbin{\oplus}}\newcommand{\xor}{\mathbin{\oplus}}
\newcommand{\F}{\mathbb F}\newcommand{\B}{\mathcal B}
\lstdefinestyle{pythonstyle}{language=Python,basicstyle=\ttfamily\footnotesize,keywordstyle=\color{Ink}\bfseries,commentstyle=\color{Muted},stringstyle=\color{Coral},showstringspaces=false,breaklines=true,frame=single,rulecolor=\color{Muted!30},columns=fullflexible}
\lstdefinestyle{cppstyle}{style=pythonstyle,language=C++}
\lstset{style=pythonstyle}
\makeindex

% Presentation layer; mathematical macros above remain unchanged.
% Shared presentation for the Quantum and Circuits monograph.
% The text block and mathematical font system are intentionally retained.
\usepackage{setspace,caption,etoolbox}
\usepackage[titles]{tocloft}
\definecolor{CoverNavy}{HTML}{142D3B}
\definecolor{Ink}{HTML}{142D3B}
\definecolor{BodyInk}{HTML}{22282C}
\definecolor{Teal}{HTML}{286B6D}
\definecolor{Muted}{HTML}{55616A}
\definecolor{Pale}{HTML}{F1F5F4}
\definecolor{Gold}{HTML}{A97D4E}
\definecolor{CoverMist}{HTML}{D7E5E5}
\colorlet{ink}{BodyInk}\colorlet{teal}{Teal}\colorlet{accent}{Teal}
\colorlet{muted}{Muted}\colorlet{pale}{Pale}\colorlet{light}{Pale}
\colorlet{deepblue}{Ink}\colorlet{restrictedblue}{Ink}\colorlet{gold}{Gold}
\colorlet{proofgreen}{Pale}\colorlet{recordblue}{Pale}\colorlet{provedgreen}{Pale}\colorlet{computedblue}{Pale}
\color{BodyInk}\setstretch{1.04}
% Reserve room for the running head while leaving the body origin unchanged.
\setlength{\headheight}{21pt}\addtolength{\topmargin}{-5pt}
\setlength{\parindent}{1.05em}
\setlength{\parskip}{2.5pt plus .5pt minus .3pt}
\setlength{\abovedisplayskip}{9pt plus 2pt minus 2pt}
\setlength{\belowdisplayskip}{9pt plus 2pt minus 2pt}
\setlength{\abovedisplayshortskip}{5pt plus 2pt minus 1pt}
\setlength{\belowdisplayshortskip}{6pt plus 2pt minus 1pt}
\clubpenalty=10000\widowpenalty=10000\displaywidowpenalty=10000\brokenpenalty=10000
\frenchspacing
% Full-width source tables start at the text-block edge, not a paragraph indent.
\BeforeBeginEnvironment{tabularx}{\noindent}
% Display math resumes horizontal mode before endproof, so the break must
% be prevented when the display is built, rather than at the later QED.
\AtBeginEnvironment{proof}{\postdisplaypenalty=10000\relax}

% A few short concluding treatments benefit from a slightly tighter rhythm.
% These scopes use chapter identifiers, never physical page numbers.
\newcommand{\VolumeChapterComposition}{%
  \setstretch{1.04}%
  \setlength{\parskip}{2.5pt plus .5pt minus .3pt}%
  \setlength{\abovedisplayskip}{9pt plus 2pt minus 2pt}%
  \setlength{\belowdisplayskip}{9pt plus 2pt minus 2pt}%
  \edef\VolumeChapterKey{\thechapter}%
  \ifdefstring{\VolumeChapterKey}{20}{%
    \setstretch{1.0}\setlength{\parskip}{1.4pt plus .4pt minus .2pt}%
    \setlength{\abovedisplayskip}{7pt plus 1pt minus 1pt}%
    \setlength{\belowdisplayskip}{7pt plus 1pt minus 1pt}%
  }{}%
  \ifdefstring{\VolumeChapterKey}{E}{%
    \setstretch{1.0}\setlength{\parskip}{1.5pt plus .4pt minus .2pt}%
    \setlength{\jot}{0pt}%
    \setlength{\abovedisplayskip}{6pt plus 1pt minus 1pt}%
    \setlength{\belowdisplayskip}{6pt plus 1pt minus 1pt}%
    \setlength{\LTpre}{5pt}\setlength{\LTpost}{5pt}%
  }{}%
}
\AtBeginEnvironment{longtable}{%
  \edef\VolumeTableChapter{\thechapter}%
  \ifdefstring{\VolumeTableChapter}{B}{\small\setstretch{1.0}}{}%
}

\newcommand{\VolumeLabel}[1]{{\sffamily\fontsize{8.5}{11}\selectfont\addfontfeatures{LetterSpace=14}\MakeUppercase{#1}}}
\newcommand{\VolumeGoldRule}{\noindent\textcolor{Gold}{\rule{42pt}{.8pt}}\par}
\newcommand{\VolumeRunningPart}{Front matter}
\fancyhf{}
\fancyhead[L]{\sffamily\fontsize{8}{10}\selectfont\color{Muted}Quantum and Circuits}
\fancyhead[R]{\sffamily\fontsize{8}{10}\selectfont\color{Teal}\MakeUppercase{\VolumeRunningPart}}
\fancyfoot[R]{\normalfont\small\color{Muted}\thepage}
\renewcommand{\headrulewidth}{.25pt}
\renewcommand{\headrule}{\hbox to\headwidth{\color{Muted!25}\leaders\hrule height \headrulewidth\hfill}}
\fancypagestyle{plain}{\fancyhf{}\fancyfoot[R]{\normalfont\small\color{Muted}\thepage}\renewcommand{\headrulewidth}{0pt}}

\titleformat{\chapter}[display]{\normalfont\color{Ink}\raggedright\hyphenpenalty=10000}{\color{Teal}\VolumeLabel{\chaptertitlename\enspace\thechapter}}{11pt}{\fontsize{25}{31}\selectfont\bfseries}
\titlespacing*{\chapter}{0pt}{5pt}{22pt}
\titleformat{\section}{\normalfont\bfseries\fontsize{14}{18}\selectfont\color{Ink}\raggedright}{\textcolor{Teal}{\thesection}}{.65em}{}
\titleformat{\subsection}{\normalfont\bfseries\fontsize{11.8}{15}\selectfont\color{Ink}\raggedright}{\thesubsection}{.65em}{}
\titleformat{\subsubsection}{\normalfont\itshape\fontsize{11.3}{14}\selectfont\color{Ink}\raggedright}{}{0pt}{}
\titlespacing*{\section}{0pt}{18pt plus 3pt minus 2pt}{8pt}
\titlespacing*{\subsection}{0pt}{13pt plus 2pt minus 1pt}{6pt}
\titlespacing*{\subsubsection}{0pt}{11pt}{4pt}

\captionsetup{format=hang,justification=raggedright,singlelinecheck=false,font={small,sf,stretch=1.04},labelfont={bf,sf,color=BodyInk},textfont={color=Muted},labelsep=period,skip=8pt}
\hypersetup{linkcolor=Teal,citecolor=Teal,urlcolor=Teal,breaklinks=true,linktoc=page}
\tcbset{enhanced,breakable,sharp corners,frame hidden,boxrule=0pt,borderline west={.8pt}{0pt}{Teal},left=11pt,right=10pt,top=9pt,bottom=8pt,before skip=11pt,after skip=11pt,fonttitle=\sffamily\bfseries\small,coltitle=Teal,colbacktitle=Pale,attach title to upper={\par\smallskip},title filled=false}

% An uncluttered table of contents with restrained part headings.
\renewcommand{\cftpartfont}{\normalfont\bfseries\color{Teal}}
\renewcommand{\cftpartpagefont}{\normalfont\color{Teal}}
\renewcommand{\cftchapfont}{\normalfont}
\renewcommand{\cftchappagefont}{\normalfont}
\renewcommand{\cftchapleader}{\hfill}
\setlength{\cftbeforepartskip}{17pt}
\setlength{\cftbeforechapskip}{5pt}
\setlength{\cftchapindent}{12pt}
\setlength{\cftchapnumwidth}{25pt}
\setlength{\cftpartnumwidth}{32pt}
\cftsetpnumwidth{2em}\cftsetrmarg{3em}

% Stable chapter anchors support the navigational lists on the part openings.
\makeatletter
\apptocmd{\ttl@mkchap@i}{\VolumeChapterComposition\ifttl@label\if@mainmatter\label{vol-ch-\thechapter}\fi\fi}{}{}
\makeatother
\newcommand{\VolumePartEntry}[2]{%
  \noindent\begin{tabularx}{\linewidth}{@{}Xr@{}}%
  #2 & \textcolor{Teal}{\pageref{vol-ch-#1}}\\%
  \end{tabularx}\par\vspace{5pt}%
}
\newcommand{\VolumePart}[3]{%
  \clearpage\refstepcounter{part}%
  \addcontentsline{toc}{part}{\protect\numberline{\thepart}#1}%
  \gdef\VolumeRunningPart{Part \thepart}%
  \thispagestyle{plain}\begingroup\setlength{\parindent}{0pt}%
  \vspace*{12pt}%
  {\color{Teal}\VolumeLabel{Part \thepart}\par}%
  \vspace{16pt}%
  {\fontsize{31}{38}\selectfont\color{Ink}\raggedright #1\par}%
  \vspace{14pt}\VolumeGoldRule\vspace{18pt}%
  {\fontsize{11.5}{16}\selectfont\color{Muted}#2\par}%
  \vspace{24pt}%
  {\color{Teal}\VolumeLabel{In this part}\par}%
  \vspace{13pt}%
  {\fontsize{10.8}{15}\selectfont #3\par}%
  \endgroup\clearpage%
}

% The first visual-guide section shares its chapter opening; later guide
% sections keep their intentional page breaks. Generated source stays intact.
\newcommand{\VolumeVisualInput}[1]{%
  \begingroup\let\VolumeSavedClearpage\clearpage%
  \def\clearpage{\let\clearpage\VolumeSavedClearpage}%
  \input{#1}\endgroup%
}

\newcommand{\VolumeCover}{%
  \begin{titlepage}\thispagestyle{empty}%
  \pagecolor{CoverNavy}\color{white}%
  \setlength{\parindent}{0pt}\setlength{\parskip}{0pt}%
  \vspace*{3pt}%
  {\color{CoverMist}\VolumeLabel{Realization Calculus}\par}%
  \vspace{22pt}\VolumeGoldRule%
  \vspace{30pt}%
  {\fontsize{52}{59}\selectfont\raggedright Quantum\\and Circuits\par}%
  \vspace{25pt}%
  {\fontsize{22}{29}\selectfont\raggedright Arithmetic, Symmetry,\\and the Cost of Realization\par}%
  \vspace{33pt}%
  {\color{CoverMist}\sffamily\fontsize{8}{11}\selectfont\addfontfeatures{LetterSpace=12}OPERATORS\enspace /\enspace GATES\enspace /\enspace OBSTRUCTIONS\par}%
  \vfill%
  {\color{white!14!CoverNavy}\rule{\linewidth}{.4pt}\par}%
  \vspace{16pt}%
  {\fontsize{12}{16}\selectfont Komorebi P.\par}%
  \vspace{7pt}%
  \end{titlepage}\nopagecolor\color{BodyInk}%
}

